\documentclass[11pt]{article}
\usepackage[T1]{fontenc}
\usepackage[utf8]{inputenc}
\usepackage{lmodern}
\usepackage{cmap}
\usepackage[a4paper,margin=2.15cm]{geometry}
\usepackage{amsmath,amssymb}
\pagestyle{plain}
\setlength{\parskip}{2.5pt}
\newcommand{\A}{\mathcal A}
\newcommand{\G}{\mathcal G}
\newcommand{\Hil}{\mathcal H}

\begin{document}

\begin{center}
{\Large\bfseries Erratum (v3): Theorems 1.1, 3.1 and 4.1\\
are not established}\\[5pt]
{\large The Yang--Mills Mass Gap on the Lattice:\\
A Conditional Synthesis}\\[6pt]
Lluis Eriksson \quad\textbullet\quad ai.viXra:2602.0033 \quad\textbullet\quad July 2026
\end{center}

\noindent\rule{\textwidth}{0.8pt}

\vspace{2pt}
\noindent
\textbf{Neither of the two pillars of this paper is established: the terminal gap
(Theorem 3.1) nor the doubling cascade that transports it (Theorem 4.1).} Theorem 1.1
is therefore \textbf{not established at several distinct points audited below}.
\textbf{What the counterexamples of this note refute are general implications used
by the proofs, one level below the numbered results.} There are \emph{two} such
implications, with different premises and different conclusions, and
\textbf{each family satisfies the abstract hypotheses relevant to the implication it
targets} --- not the union of both hypothesis sets.

\textbf{The exact-trace family targets the exact/additive extraction} (\S2):
positivity, self-adjointness, trace class, and the trace identity holding
\emph{exactly} for infinitely many $M$, hence $\bar\delta=0$ and the printed window
$M_0\le M$, $e^{-m_nM}\gg\bar\delta$ of the displayed estimate --- the manuscript
prints $\gg$ there, and with $\bar\delta=0$ it holds however it is read. It does \textbf{not}
satisfy the cap $m_nM\le\kappa'/2$ of part (i) --- its gap is a fixed $a>0$, so
$m_nM\to\infty$ --- and it does not need to: what it answers is the uncapped
exact/additive step.
\textbf{The finite-error family targets the \mbox{finite-window} relative extraction} (\S3):
the approximate identity within the admitted error floor, the printed lower bound on
the window, and the printed one-sided cap, for every $\kappa'\ge1$. Each contradicts
the conclusion of the implication it is aimed at, so \textbf{both abstract extraction
implications are false}. \emph{Singular --- ``the abstract extraction implication'' ---
survived into the round that separated the two, because it had been made a required
phrase the round before; a guard can preserve an error as faithfully as it preserves a
correction.} \emph{An earlier draft of this note said both
families met every printed hypothesis. They do not, and they need not: a
counterexample answers the implication it is aimed at, and stating more than that
overstates it in the direction this whole correction campaign exists to avoid.}

It does not follow that a numbered result is false: Theorem 4.1 and Proposition 4.4
are stated about the operators $T_n$ built in Step 1 from Balaban's effective actions,
neither is stated for an arbitrary positive trace-class family, and \textbf{none is
shown to arise from the Balaban effective actions} or to coincide with the intended
operators. \emph{Nor is the converse claimed: that no such family can so arise is not
established here either.} Accordingly, \textbf{the conclusion for the intended
Yang\nobreakdash-Mills transfer operators is not established, not refuted here.} \emph{A previous round of this note
declared the conclusion of Proposition 4.4(i) refuted outright, on this page and in
the submission metadata. That went one level too far and is withdrawn; what changed is
the classification, not the mathematics, and the families themselves stand.} The exact
status, piece by piece:
\begin{itemize}\setlength{\itemsep}{0pt}
\item \textbf{the exact/additive extraction implication} used in Steps 3--4 of
Theorem 4.1 and again in the display of Proposition 4.4 --- \textbf{false}, by the
exact-trace family;
\item \textbf{the \mbox{finite-window} relative extraction implication} used in Proposition
4.4(i) --- \textbf{false} under the standard uniform reading, by the finite-error
family. \emph{These two are defects of one architecture, but they are not one
implication: their premises differ and so do their conclusions};
\item the additive estimate displayed in Proposition 4.4 --- not established by the
printed proof, and no survivor either: \textbf{the exact-trace family also contradicts
the displayed additive estimate} read as a general implication;
\item Proposition 4.4(i), the relative cascade --- not established; as a general
implication it fails quantitatively and not merely as an exact equality, the second
family keeping $\dfrac{m_{n+1}}{2m_n}-1=1$ rather than $O(e^{-\kappa'/2})$;
\item Proposition 4.4(ii), the extensive regime --- the two families do not reach its
volume regime; it is \emph{likewise not established by the printed proof}, since it
rests on the same additive extraction whose derivation is invalid through the
volume-dependent operator. Its separate volume-cap algebra is not audited here;
\item \textbf{(H-FACT) as consumed by Proposition 4.4} --- the printed prefactor
$C_{\rm fact}$ is normalised to one; for $C_{\rm fact}>1$ that needs
$\kappa'>\log C_{\rm fact}$ and $\kappa'-\log C_{\rm fact}$ everywhere downstream, and
no such renormalisation is stated. \textbf{Nor does the printed hypothesis specify
coupling uniformity} (\S3);
\item Lemma 4.3(a) --- a conditional algebraic identity;
\item Lemma 4.3(b) --- an invalid inference, using a lower bound as an upper bound.
\end{itemize}
\emph{Nothing here claims the physical gap differs from what the paper expected.}
\textbf{For Theorems 1.1, 3.1, 4.1 and Proposition 4.4, what is withdrawn here is the
printed derivation; no conclusion of theirs about the intended operators is refuted.}
\emph{What is false here, classified rather than counted: the exact/additive
extraction implication; the \mbox{finite-window} relative extraction implication; and the
printed orbit-space dimension corrected in \S1.3. None of the three is a conclusion of
a numbered result about the intended operators, and this note does not claim the list
is closed --- a closing count is the failure this chain has now recorded twice.}
\emph{This audit does not
refute} (H-BAL) \emph{or} (H-FACT), but those hypotheses do not suffice for the
printed conclusions without the additional operator and identification inputs stated
below.

\emph{An earlier draft of this note stated that ``the doubling cascade of \S4 is
unaffected and is the part of the paper that survives intact''. That is withdrawn:
\S2 below shows it is not.} The provenance page appended here fixes the dictionary
that \S1.1 and \S4.1 turn on.

\noindent\textbf{Notation, before anything else.} \textbf{In this erratum $\nu_{\rm push}:=\pi_\#\mathrm{Vol}_{\A}$, and every $\nu$ below means $\nu_{\rm push}$.} The restored manuscript uses the same letter for a different object: Step 3 of Theorem 3.1 calls $\nu_{\rm Vol}:=d\mathrm{Vol}_B/\mathrm{Vol}_B(B)$ ``the uniform measure $\nu$''. \textbf{The two are not identified}, and identifying them is what (I2b) of the appended page forbids. \emph{Read one for the other and a printed derivation looks as though it already used the corrected measure.}

\noindent\textbf{And four different constants are called $C$ in this chain.} The
manuscript writes $C(N_c)=24\pi^2/(11N_c)$ for the exponent of Theorem 1.1, $C$ for the
prefactor of (H-FACT), $C$ again for the lower end of the window in Proposition 4.4
(``$C\le M$''), and $C_0$ for the volume constant of Remark 2.2. \textbf{This erratum
writes $C_{\rm fact}$ for the (H-FACT) prefactor, $M_0$ for the window's lower end, and
$C_\beta$ for the terminal-coupling constant it introduces in \S4.2}, keeps $C(N_c)$ and
$C_0$ as printed, and writes a bare $C$ only here and when quoting the manuscript's own
``$C\le M$''. \emph{An earlier draft of this note promised never to write a bare $C$ and
then wrote one, in \S4.2, in the very round that put up the barrier.}

\section{Theorem 3.1: the terminal gap}

\subsection{The printed Euclidean measure is not identified with the ground-state
measure}

Step 2 introduces $d\mu_{\rm print}=e^{-f}d\mathrm{Vol}_B/Z$ with
$f=\beta_{n_{\max}}V_W+\sum_X\epsilon_{n_{\max}}(X)$ --- the \emph{effective action} ---
applies Holley--Stroock in Step 3, and writes
$m_{n_{\max}}=\frac{g_{n_{\max}}^2}{2}\lambda_1$. That relation belongs to the
Hamiltonian framework, where the ground-state transformation requires the measure to
carry $f_0=-\log\psi_0^2$, i.e.\ $\mu_0=\psi_0^2\nu_{\rm push}$.

The Gibbs measure of a Euclidean effective action is \emph{a priori} a distinct
object, and its weighted Dirichlet-form gap \emph{a priori} a distinct spectral
quantity. \textbf{The printed derivation applies the Hamiltonian gap identity to an
effective-action Gibbs measure without identifying that measure with the ground-state
measure}: no identification of $\mu_{\rm print}$ with $\mu_0=\psi_0^2\nu_{\rm push}$ has been
proved. In special models they may coincide, but in general the passage requires a
transfer-operator or Osterwalder--Seiler identification theorem, not a change of
notation. This is (I3) of the appended page.

\emph{Note the reference measure.} What Step 2 prints is against
$\mathrm{Vol}_B$, not against $\nu_{\rm push}$, and it is a \emph{regular-stratum}
formula: no
global extension, and no closed Dirichlet form for it, is specified. On $B_{\rm reg}$,
$d\nu_{\rm push}=w\,d\mathrm{Vol}_B$ with $w$ the orbit-volume density, so
\[
d\mu_{\rm print}=Z_{\rm print}^{-1}e^{-S_{\rm eff}}d\mathrm{Vol}_B
=Z_{\rm print}^{-1}e^{-(S_{\rm eff}+\log w)}d\nu_{\rm push}
\qquad\text{on }B_{\rm reg},
\]
which is an exact rewriting of one measure, with the same normaliser. Relative to
$\nu_{\rm push}$ the printed measure therefore carries the exponent
$S_{\rm eff}+\log w$, while the candidate $\nu_{\rm push}$-based measure
$d\widetilde\mu=\widetilde Z^{-1}e^{-S_{\rm eff}}d\nu_{\rm push}$ carries
$S_{\rm eff}$; they
coincide only under an additional constancy or identification statement for $w$, none
of which is established here. This matters for the repair, and is (I2b) of the
appended page.

\emph{The programme already knew.} The v2 of \texttt{ai.viXra:2602.0035} corrected
exactly this substitution in its own v1 --- the action-based measure being
``\emph{dimensionally inconsistent}'' and, per its suite, failing to reproduce the gap
at ``\emph{$>50\%$, growing}''. That correction did not reach this paper, though both
are July 2026 documents of the same series. \emph{What that measurement is, and is
not:} it concerns the literal potential-based measure of that paper's v1, not the
$\epsilon$-corrected effective-action measure printed here. \textbf{It supplies strong numerical counterevidence to a
universal action-measure substitution, and reinforces that no identification
may be made by notation; it is neither a proof that no general identification
exists nor an analytic counterexample to Theorem 3.1 of the present paper.}
\emph{An earlier draft said it ``confirms'' the unavailability of a general
identification --- a non-rigorous one-plaquette computation cannot confirm a
mathematical non-existence. The reason to withdraw the printed derivation is
that no identification theorem has been proved, not that a number came out
wrong.}

\subsection{The geometric input is no longer available}

Steps 1 and 3 use $\operatorname{Ric}_B\ge N_c/4$ from
\texttt{ai.viXra:2602.0036}, Theorem 3.1. That theorem is stated \emph{at regular
points} while Bakry--\'Emery is applied here globally, and the corollary that
globalised it has been withdrawn. Decisively, \textbf{that theorem is itself not
established by its printed proof}: its O'Neill trace omits the mixed
horizontal--vertical sectional terms, which for a bi-invariant metric on a compact
group are non-negative and subtract. \textbf{No pointwise Ricci bound on the orbit
space may presently be cited.}

What is available is \texttt{ai.viXra:2602.0046} v3: a global curvature-dimension
bound, hence a Poincar\'e/log-Sobolev inequality, for $(B,d_B,\nu)$ with
$\nu=\pi_\#\mathrm{Vol}_{\A}$ --- a bound for $K_\nu$, not for a realisation of the
unweighted $\mathcal K_B^{\rm reg}$. It cannot serve as the pointwise input used here.

\emph{A warning about the symbol $\nu$, and it is load-bearing.} \textbf{The $\nu$ of
the appended page is not the $\nu$ of the restored pages.} That page writes
$\nu_{\rm push}=\pi_\#\mathrm{Vol}_{\A}$; Step 3 of Theorem 3.1 calls
$\nu_{\rm Vol}=d\mathrm{Vol}_B/\mathrm{Vol}_B(B)$ ``the uniform measure $\nu$''.
Reading one for the other makes the printed derivation look as though it already used
the corrected measure --- which is exactly the identification (I2b) forbids.
\emph{Nor does it repair Step 3 by a change of reference measure}: by (I2b),
\textbf{replacing $d\mathrm{Vol}_B$ by $d\nu_{\rm push}$ under a fixed exponent is
not a free
rewriting}. Two routes are open and they are different theorems.
\textbf{(A)} Keep $\mu_{\rm print}$, whose exponent relative to $\nu_{\rm push}$ is
$S_{\rm eff}+\log w$: one must then extend the regular-stratum formula globally and
bound $\operatorname{osc}(S_{\rm eff}+\log w)$, for which the bound on
$\operatorname{osc}(S_{\rm eff})$ does not suffice --- if $w$ degenerates at strata of
enlarged stabiliser, $\log w$ is exactly the term that destroys a bounded-perturbation
argument, and \textbf{no control of $\log w$ is supplied} in this chain.
\textbf{(B)} Reformulate over $(B,\nu_{\rm push})$ with
$d\widetilde\mu=\widetilde Z^{-1}e^{-S_{\rm eff}}d\nu_{\rm push}$, where \texttt{2602.0046} does
supply the base inequality --- but that is not the measure the paper printed, and a
theorem is then owed showing it is the correct Euclidean object.

\subsection{A four-dimensional action on a three-dimensional orbit space}

Step 1 fixes the space as $B_{n_{\max}}$ and prints its dimension as
$3L_{n_{\max}}^3(N_c^2-1)$. \textbf{That is the dimension of the spatial
configuration manifold} $\A_{\rm spat}=SU(N_c)^{3L^3}$, \textbf{not of the orbit
space}: the site gauge group $SU(N_c)^{L^3}$ has dimension $L^3(N_c^2-1)$ and acts
with discrete generic stabiliser, so the \textbf{principal-stratum dimension is
$2L_{n_{\max}}^3(N_c^2-1)$}, and at reducible configurations the quotient is
stratified with no single manifold dimension. \emph{An earlier draft of this note
repeated the printed figure as though it were correct.} The $O(L^3)$ versus
$O(L^4)$ mismatch below is unchanged by the correction. Step 1 also says
``\emph{the spatial slice is three-dimensional}''; Step 2 bounds the oscillation with
$O(L_{n_{\max}}^4)$ plaquettes. Both are $O(1)$ terminally, so no numerical bound
breaks --- but the mismatch shows $f$ is the four-dimensional Euclidean effective
action while $\mu$ is treated as living on the three-dimensional spatial orbit space.
(H-FACT) approximately controls the factorisation of the \emph{partition function}; it
does not build a Hamiltonian on $B^{\rm spat}$ nor identify its ground state.

\section{Theorem 4.1: the limit is taken over operators that change with the volume}

This is the defect that removes the last standing pillar, and it is independent of
everything in \S1.

Step 1 defines $T_n$ on $\Hil_n=L^2(B_n^{\rm spat})$ for
$\Lambda_n=(\mathbb Z/L_n\mathbb Z)^4$ with $L_n=L/2^n$, using --- in the paper's own
words --- ``\emph{the convention of equal spatial and temporal extent}''. \textbf{Both
the Hilbert space and the operator therefore depend on $L_n$.}

Step 2 sets $M:=L_n$ and obtains
$\operatorname{Tr}(T_{n+1}^{M/2})=\operatorname{Tr}(T_n^{M})$, then justifies
$M\to\infty$ by \emph{varying the original isotropic size $L$}. But varying $L$ varies
$L_n=M$, hence varies the spatial slice, hence varies $\Hil_{n}$, $T_{n}$, the
eigenvalues $\lambda_j^{(n)}$ and the gap $m_n$ --- all of which Step 3 then treats as
constants in $M$ when it writes
\[
\operatorname{Tr}(T_n^M)=(\lambda_0^{(n)})^M\bigl(1+e^{-m_nM}+O(e^{-m_n'M})\bigr).
\]
That expansion extracts the spectrum only when a \emph{fixed} operator is raised to a
growing power. Here the power and the operator grow together.

\paragraph{The trace identities do not imply the conclusion.} An abstract example
suffices to show that \textbf{the doubling and the displayed additive estimate do not
follow from the family of identities alone} --- not that the identities carry no
spectral information at all, which would be a stronger claim and is not made here. For \emph{fixed} $a,b>0$ with $b\neq 2a$, and for each even $M$, put
\[
T_{n,M}=\begin{pmatrix}1&0\\0&e^{-a}\end{pmatrix},\qquad
T_{n+1,M/2}=c_M\begin{pmatrix}1&0\\0&e^{-b}\end{pmatrix},\qquad
c_M=\Bigl(\frac{1+e^{-aM}}{1+e^{-bM/2}}\Bigr)^{2/M}.
\]
Both matrices are positive, self-adjoint and trace class; \emph{positivity of $a$ and
$b$ is what makes the top eigenvalue the first entry}, so the gaps are exactly $a$ and
$b$. Then $\operatorname{Tr}(T_{n+1,M/2}^{M/2})=\operatorname{Tr}(T_{n,M}^{M})$ for
\emph{every} even $M$, while $b$ ranges over all positive values other than $2a$. This
is not offered as a model of anything physical, and \textbf{it is not shown to arise
from any lattice gauge theory} --- nor is the converse shown; what it does show is that
the inference used in Steps 3--4 is not valid for an $M$-dependent family.

\paragraph{With a third eigenvalue, and against the additive estimate as well.} The
two-by-two family leaves the printed remainder $O(M^{-1}e^{-(m_n'-m_n)M})$ with no
third eigenvalue to refer to, which invites the objection that it is testing a
degenerate case. It is not needed. Fix $a,b,r,s>0$ with $b\neq2a$, $r>a$ and $s>b$,
and for each even $M$ put
\[
T_{n,M}=\mathrm{diag}(1,e^{-a},e^{-r}),\qquad
T_{n+1,M/2}=c_M\,\mathrm{diag}(1,e^{-b},e^{-s}),
\]
\[
c_M=\Bigl(\frac{1+e^{-aM}+e^{-rM}}{1+e^{-bM/2}+e^{-sM/2}}\Bigr)^{2/M}.
\]
Then $\operatorname{Tr}(T_{n+1,M/2}^{M/2})=\operatorname{Tr}(T_{n,M}^{M})$
\emph{exactly}, for every even $M$, so one may take $\delta_n=\delta_{n+1}=0$; the
orderings $1>e^{-a}>e^{-r}$ and $1>e^{-b}>e^{-s}$ hold, and the gaps and next distinct
gaps are $m_n=a$, $m_{n+1}=b$, $m_n'=r$, $m_{n+1}'=s$. The displayed additive estimate
of Proposition 4.4 would give
\[
b=2a+O\bigl(M^{-1}e^{-(r-a)M}\bigr),
\]
whose right-hand side tends to $2a$ along the infinitely many admissible $M$,
contradicting $b\neq2a$. \textbf{So the exact-trace family also contradicts the
displayed additive estimate}, and not only the relative cascade built on it.
\emph{Verified at $a=1/3$, $b=1$, $r=2$, $s=5$ for $M=4,8,20,60,200,1000$: the two
traces agree to sixty digits while $|b-2a|=1/3$ stays fixed and the printed remainder
falls below $10^{-700}$.} \emph{An earlier draft of this note recorded the additive
estimate as untouched by these families; \S3 said so too.}

\paragraph{What these families do and do not refute.} They are matrices, not the
transfer operators of Step 1: \textbf{none is shown to arise from the Balaban effective
actions}, and this note does not claim they arise from any lattice gauge theory --- nor
does it claim they cannot. \textbf{Each family satisfies the abstract hypotheses
relevant to the implication it targets}, and only those: the exact-trace family has a
fixed gap $a>0$, so $m_nM\to\infty$ and the one-sided cap of part (i) is not among its
hypotheses, while the finite-error family of \S3 does satisfy that cap. What they refute are two \emph{general} implications.
\textbf{The exact-trace family targets the exact/additive extraction}: positivity,
self-adjointness, trace class and the exact trace identity for infinitely many $M$,
therefore the doubling and the displayed additive estimate --- which is what Steps
3--4 use. \textbf{The finite-error family targets the
\mbox{finite-window} relative extraction}: the approximate identity, the error floor and the one-sided window,
therefore a relative cascade with error $O(e^{-\kappa'/2})$ --- which is what
Proposition 4.4(i) uses. Theorem 4.1 and Proposition 4.4 are stated about the
operators built in Step 1, not about that abstract class, so their conclusions are
\emph{not established} for the intended operators rather than refuted. \emph{This
distinction is the one the note previously collapsed, calling the conclusion of
Proposition 4.4(i) false outright.}

\paragraph{A second, local defect of Step 3: multiplicities.} Even for a fixed
operator the printed expansion is not correct as written. It sets the coefficient of
$e^{-m_nM}$ to $1$ and takes $m_n'=-\log(\lambda_2/\lambda_0)>m_n$, then adds that if
$\lambda_1$ is degenerate the coefficient should be its multiplicity. But with the
printed enumeration a degenerate $\lambda_1$ gives $\lambda_2=\lambda_1$ and hence
$m_n'=m_n$, not $m_n'>m_n$: the remainder is then of the same order as the term it is
supposed to be smaller than. Nor does ``positive, self-adjoint and trace class'' make
$\lambda_0$ simple --- that needs a positivity-improving or Perron--Frobenius input.
The correct form groups \emph{distinct} eigenvalues $\lambda_0>\lambda_1>\lambda_2$
with multiplicities $d_j$:
\[
\operatorname{Tr}(T^M)=d_0\lambda_0^M\Bigl[1+\tfrac{d_1}{d_0}e^{-mM}
+O(e^{-m'M})\Bigr],\qquad m'>m .
\]
This does not change the verdict --- Theorem 4.1 is already not established by the
$M$-dependence above --- but it is recorded because this note claims to have audited
the cascade, and because the corrected form is what any repair must use.

\paragraph{Why the suite does not see it.} The verification miniature is a
one-dimensional $\mathbb Z_N$ clock chain under decimation. In one dimension, varying
the length varies only the temporal direction: there is no growing spatial slice, so
the operator stays fixed and the extraction is legitimate. The miniature therefore
tests a situation from which the defect is absent by construction.

\paragraph{What would repair it.} \textbf{A fixed spatial slice by itself is not
sufficient.} A family of \emph{rectangular} lattices with the spatial slice held fixed
and only the temporal extent varied removes the $M$-dependence of $\Hil_n$ and $T_n$,
which is the defect above; but the extraction still needs an identity to feed it, and
the printed one is $\operatorname{Tr}(T_{n+1}^{M/2})=\operatorname{Tr}(T_n^{M})$
obtained by varying the \emph{isotropic} size. \emph{Nor is a fixed slice
necessary: the last two routes below do not presuppose one, and an earlier draft of this
note asserted a necessity it had not shown.} Three routes, then:
\begin{itemize}\setlength{\itemsep}{0pt}
\item \textbf{a fixed-spatial-slice family \emph{together with} an RG/transfer-matrix
identity valid for all temporal extents over that fixed space} --- both, not either;
\item or an operator relation such as $T_{n+1}$ conjugate to $T_n^2$, which gives the
doubling directly and needs no limit;
\item or \textbf{a quantitative, normalised spectral-convergence theorem} for the
volume-dependent family. \textbf{Ordinary spectral convergence is insufficient}, and
the exact-trace family above already refutes it: there $\log c_M\to0$, so
$T_{n+1,M/2}\to\mathrm{diag}(1,e^{-b},e^{-s})$ in operator norm while $T_{n,M}$ does not
move at all; every eigenvalue converges, the limiting gaps are $a$ and $b$, and still
$b\neq2a$. The reason is that $c_M-1\sim(2/M)\bigl(e^{-aM}-e^{-bM/2}\bigr)$ is of the
\emph{same} exponential order as the first excited contribution, and absorbs it exactly.
What a repair needs is
\textbf{leading-eigenvalue normalisation relative to the actual finite-$M$
first-excited contributions}, \emph{and it must carry the multiplicities this
note has just insisted on}. Write $d_{0,n}(M)$ for the multiplicity of the top
eigenvalue, so that the leading trace term is $d_{0,n}(M)\lambda_{0,n}(M)^M$, and set
\[
 \begin{aligned}
 Q_M&:=\frac{d_{0,n+1}(M/2)\,\lambda_{0,n+1}(M/2)^{M/2}}
 {d_{0,n}(M)\,\lambda_{0,n}(M)^{M}},\\
 A_n(M)&:=\frac{d_{1,n}(M)}{d_{0,n}(M)},\qquad
 A_{n+1}(M/2):=\frac{d_{1,n+1}(M/2)}{d_{0,n+1}(M/2)} .
 \end{aligned}
\]
Suppose $m_n(M)\to m_n^\infty>0$ and
$m_{n+1}(M/2)\to m_{n+1}^\infty>0$.  The normalisation condition required by the
trace comparison is
\[
 \boxed{\ |Q_M-1|=o\!\left(
 A_n(M)e^{-m_n(M)M}
 +A_{n+1}(M/2)e^{-m_{n+1}(M/2)M/2}\right).\ }
\]
The earlier condition $Q_M=1+o(e^{-\mu_\infty M})$, with
$\mu_\infty=\min\{m_n^\infty,m_{n+1}^\infty/2\}$, is \textbf{necessary-looking but
not sufficient}: it compares with a limiting exponential scale rather than the
finite-$M$ corrections the trace must resolve.  The boxed condition is required
together with convergence of the ratios of \emph{distinct} eigenvalues,
\textbf{subexponential control of the excited multiplicities}
\[
\tfrac1M\bigl|\log\bigl(d_{1,n}(M)/d_{0,n}(M)\bigr)\bigr|\to0,
\qquad
\tfrac1M\bigl|\log\bigl(d_{1,n+1}(M/2)/d_{0,n+1}(M/2)\bigr)\bigr|\to0,
\]
and \textbf{uniform control of the total excited trace tail}
\[
\sum_{j\ge2}\frac{d_{j,n}(M)}{d_{1,n}(M)}
\,e^{-[m_{j,n}(M)-m_{1,n}(M)]M}\longrightarrow0 ,
\]
and the same at the next scale with $M/2$ in place of $M$ throughout.
For an exact trace identity these conditions do force the limiting doubling.  Indeed,
put
\[
 X_M=A_n(M)e^{-m_n(M)M},\qquad
 Y_M=A_{n+1}(M/2)e^{-m_{n+1}(M/2)M/2}.
\]
The two tail conditions reduce the normalised trace identity to
$1+X_M(1+o(1))=Q_M[1+Y_M(1+o(1))]$.  Here both $o(1)$ terms are
\textbf{multiplicative relative-tail errors}; an additive absolute $o(1)$ would be
too weak because it could absorb either exponentially small first-excited term.  The
boxed condition then gives
$X_M-Y_M=o(X_M+Y_M)$, hence $X_M/Y_M\to1$.  Taking logarithms, dividing by $M$,
and using the two subexponential multiplicity conditions yields
$m_n^\infty=m_{n+1}^\infty/2$.  Before passing to the limit the same argument gives
the stronger exponent matching
\[
 m_n(M)M-\tfrac12m_{n+1}(M/2)M-\log\!\frac{A_n(M)}{A_{n+1}(M/2)}\longrightarrow0.
\]
This is the missing sufficiency argument; without the boxed finite-$M$ comparison,
the conclusion does not follow.
{\sloppy \textbf{\mbox{Level-by-level multiplicity} \mbox{control is not sufficient}}: it bounds each row of an enumeration, not how
many rows there are.\par}
\textbf{Positivity improving removes only the \mbox{top-eigenvalue}
multiplicity}; it says nothing about $d_1,d_2,\dots$, nor about the spectral counting
function.

\emph{The multiplicity and tail conditions are both necessary, and each is forced by
its own counterexample; they are not sufficient without the boxed normalisation.} The
first family below has one excited level of exponentially large multiplicity; the second
has exponentially many simple excited levels. \textbf{Neither the multiplicity condition
nor the tail condition excludes both.}

\emph{First: multiplicity.} For even $M$ put $D_M=\lfloor e^{M}\rfloor$, $q=e^{-2}$, and
let $P_N=\tfrac1N\mathbf1\mathbf1^{*}$. Take
\[
T_{n,M}=q\,I_{D_M+1}+(1-q)P_{D_M+1},\qquad
T_{n+1,M/2}=c_M\bigl(q\,I_2+(1-q)P_2\bigr),
\]
with $c_M=\bigl[(1+D_Me^{-2M})/(1+e^{-M})\bigr]^{2/M}$. Every entry is strictly
positive, so both are \emph{positivity improving}; the spectra are $1$ simple and
$e^{-2}$ with multiplicity $D_M$, respectively $1$ and $e^{-2}$ simple. The traces agree
exactly, both equal to $1+D_Me^{-2M}$, so $\delta_n=\delta_{n+1}=0$; both gaps are $2$,
so $m_{n+1}^\infty\neq2m_n^\infty$; $d_{0,n}=d_{0,n+1}=1$; the distinct-eigenvalue ratios
are constant; and $\mu_\infty=\min\{2,1\}=1$ while the normalisation quotient is
$c_M^{M/2}=(1+D_Me^{-2M})/(1+e^{-M})=1+O(e^{-2M})=1+o(e^{-\mu_\infty M})$.
\textbf{The normalisation condition and the positivity-improving alternative are both
satisfied, and the doubling still fails.} The mechanism is the excited multiplicity:
$D_Me^{-2M}\sim e^{-M}$, so the trace sees the effective exponent
$2-\tfrac1M\log D_M\to1$, not the spectral gap $2$. \emph{Here
$\tfrac1M\log(d_1/d_0)\to1$; there is no level $j\ge2$, so the tail sum is empty and
does not see it.}

\emph{Second: counting.} Keep $D_M=\lfloor e^{M}\rfloor$ and choose distinct
$\varepsilon_{j,M}\in[-\tfrac14,\tfrac14]$ with $\sum_{j\le D_M}\varepsilon_{j,M}=0$ ---
equispaced and symmetric will do --- and order them so that
$\varepsilon_{1,M}=\max_j\varepsilon_{j,M}$. Put
$\lambda_{j,M}=e^{-2}(1+\varepsilon_{j,M})^{1/M}$
and
\[
T_{n,M}=\mathrm{diag}\bigl(1,\lambda_{1,M},\dots,\lambda_{D_M,M}\bigr),\qquad
T_{n+1,M/2}=c_M\,\mathrm{diag}(1,e^{-2}),
\]
with the same $c_M$. Then $\operatorname{Tr}(T_{n,M}^M)=1+e^{-2M}\sum_j(1+\varepsilon_j)
=1+D_Me^{-2M}$, the traces agree exactly again, \textbf{every multiplicity is one},
the excited ratios converge to $e^{-2}$ uniformly, both limiting gaps are $2$, and the
normalisation quotient is the same $1+o(e^{-\mu_\infty M})$. \textbf{So the whole
multiplicity condition is satisfied trivially --- and the doubling still fails.} What
fails now is the tail: each term of $\sum_{j\ge2}$ lies in $[\tfrac35,1]$ and
there are $D_M-1$ of them, so the sum grows like $e^{M}$. \emph{Checked at
$M=4,8,16,24$: the traces agree, the normalisation error
$|Q_M-1|/e^{-\mu_\infty M}$ falls to $3.2\times10^{-11}$, and the tail sum exceeds
$1.6\times10^{10}$.} \emph{An earlier draft of this note wrote ``and likewise for the
levels making up the remainder'', which is not a verifiable hypothesis; the defect it
was gesturing at is the counting function, and this family is what makes that
concrete.} \textbf{Note the
arguments: the operator compared at the next scale is $T_{n+1,M/2}$, so its leading
eigenvalue is evaluated at $M/2$, and the gaps are the limits, not constants in $M$.} Or else an
identity already normalised so that the leading eigenvalue cannot absorb the excited
term.

\emph{Third: the limiting scale is not the finite-$M$ correction.}  Let
$P_2=\tfrac12\mathbf1\mathbf1^*$ and
$A(t)=e^{-t}I_2+(1-e^{-t})P_2$, whose entries are strictly positive and whose
eigenvalues are $1,e^{-t}$.  Fix $a=1/3$, $b=1$, and for even $M$ put
\[
 a_M=a+M^{-1/2},\qquad b_{M/2}=b+(M/2)^{-1/2},
\]
\[
 T_{n,M}=A(a_M),\qquad T_{n+1,M/2}=c_M A(b_{M/2}),\qquad
 c_M=\left(\frac{1+e^{-a_MM}}{1+e^{-b_{M/2}M/2}}\right)^{2/M}.
\]
The traces agree exactly; both operators are positivity improving; all top and first
excited multiplicities are one; the tail is empty; and the gaps converge to $a$ and
$b$.  With $\mu_\infty=1/3$, the old condition still holds:
$Q_M-1=O(e^{-M/3-\sqrt M})=o(e^{-\mu_\infty M})$.  Nevertheless
$b\ne2a$.  Here $Q_M-1$ is of the same order as the actual dominant finite-$M$
excited correction, so the boxed condition fails exactly where it should.  Thus
positivity improving does not rescue the R29 condition.

\emph{The original fixed-gap exact-trace family of Section 2 has $d_0=1$ at both scales and gives
$1+\Theta(e^{-\mu_\infty M})$ in that quotient, so the condition is not vacuous: it
excludes it.} \emph{An earlier draft of this note wrote the numerator at $M$ rather than
$M/2$, and the gaps as constants --- in a note whose central defect is confusing what
moves with $M$ --- and normalised only the leading eigenvalues, ignoring the very
multiplicity correction it had just declared indispensable.} \emph{An earlier draft of this note offered plain uniform
spectral convergence as a repair, which its own counterexample refutes.}
\end{itemize}
None appears. \emph{An earlier draft of this note listed the first item as one
alternative among four, which understates what it owes.}

\textbf{Theorem 4.1 is not established by its printed proof.}

\section{What this costs}

\paragraph{Proposition 4.4 silently normalises the prefactor of (H-FACT) to one.}
(H-FACT) is printed as
\[
|\delta_n|\le C_{\rm fact}\,e^{-\kappa'},\qquad
\kappa'>0\ \text{independent of }L_n .
\]
Proposition 4.4(i) invokes ``\emph{if $\bar\delta\le e^{-\kappa'}$ uniformly in $L_n$
(Hypothesis (H-FACT))}''. \textbf{The prefactor has been set to one.} The implication
from $C_{\rm fact}e^{-\kappa'}$ to $e^{-\kappa'}$ is \emph{automatic when
$C_{\rm fact}\le1$}; for $C_{\rm fact}>1$ it needs the renormalised exponent
\[
\widetilde\kappa=\kappa'-\log C_{\rm fact} ,
\]
hence \textbf{$\kappa'>\log C_{\rm fact}$} to leave a positive exponent at all, and then
$\widetilde\kappa$ in place of $\kappa'$ \emph{everywhere downstream}: the cap of
Lemma 4.3 becomes $m_nM\le\widetilde\kappa/2$ and the cascade factor becomes
$1+O(n_{\max}e^{-\widetilde\kappa/2})
=1+O(n_{\max}\sqrt{C_{\rm fact}}\,e^{-\kappa'/2})$, so $C_{\rm fact}$ enters the
accumulation as $\sqrt{C_{\rm fact}}$. \textbf{No such renormalisation is stated
anywhere}, and the manuscript nowhere bounds $C_{\rm fact}$. \emph{An earlier draft of
this note called the substitution ``a strictly stronger bound''; that is not universal
--- it is an equality of content when $C_{\rm fact}=1$ and a weakening when
$C_{\rm fact}<1$.}
\emph{This is independent of the accumulation and of the resolution defects below:
before any iteration, the proof does not literally consume the hypothesis its own
theorem declares.} The counterexample below is stated against the proposition's own,
stronger, printed bound $\bar\delta\le e^{-\kappa'}$, so it is untouched by this.

\paragraph{Proposition 4.4 --- the window controls the wrong side.} Its additive
estimate has the shape
$m_{n+1}=2m_n+O(\bar\delta e^{m_nM}/M)+O(e^{-(m_n'-m_n)M}/M)$. Writing $x=m_nM$ and
dividing by $m_n=x/M$ to get the \emph{relative} estimate the cascade then uses,
what is needed is
\[
\bar\delta\,e^{m_nM}\ll m_nM
\qquad\text{and}\qquad
e^{-(m_n'-m_n)M}\ll m_nM .
\]
The paper instead requires only
$e^{-x}\gg\bar\delta$ and then caps $x\le\kappa'/2$. \textbf{Neither of the two
conditions follows}: as $x\to0$ the printed condition becomes automatically easy,
since $e^{-x}\to1$, while $\bar\delta e^x/x$ grows without bound. \textbf{A fixed
error floor can hide a positive gap}, and a very small gap is not an advantage
against a fixed error --- it falls below the resolution.

\emph{A finite positive counterexample, as a family parametrised by $\kappa'$.}
\textbf{Put $\bar\delta=e^{-\kappa'}$ and $x=\bar\delta/2$}, and choose any
\textbf{even $M\ge M_0$}, the printed lower bound on the window (the manuscript's
$C$ in ``$C\le M$''). Take
$T_n=\mathrm{diag}(1,e^{-x/M})$ and $T_{n+1}=\mathrm{diag}(1,e^{-4x/M})$, with gaps
$m_n=x/M$ and $m_{n+1}=4x/M$, so $m_{n+1}\neq2m_n$. Put $\delta_n=0$ and
$\delta_{n+1}=(1+e^{-x})/(1+e^{-2x})-1$, which satisfies
$0<\delta_{n+1}<x<\bar\delta$. Then
$\operatorname{Tr}(T_{n+1}^{M/2})(1+\delta_{n+1})
=\operatorname{Tr}(T_{n}^{M})(1+\delta_n)$ exactly: the approximate partition
identity holds within the admitted errors, both gaps are positive, $e^{-x}\gg
\bar\delta$ holds with margin $e^{\kappa'}$ up to a factor tending to $1$, the window
$x\le\kappa'/2$ holds \textbf{for every $\kappa'\ge1$}, \emph{and the doubling
fails}.
\emph{Not for every positive value of the parameter: that window is
$e^{-\kappa'}\le\kappa'$, first true at $\kappa'=W(1)=0.5671\ldots$, and the same
condition one scale up, $2x\le\kappa'/2$, is $2e^{-\kappa'}\le\kappa'$, first true at
$W(2)=0.8526\ldots$; an earlier draft asserted it for all positive values, which is
false below those two thresholds. Taking $\kappa'\ge1$ clears both, and clears
$e^{-2x}\ge\bar\delta$ at the next scale as well.}

\emph{And it fails relatively, not merely exactly.} Proposition 4.4(i) does not
promise exact doubling: it promises $1+O(e^{-\kappa'/2})$ per step. Here
$m_n=x/M$ and $m_{n+1}=4x/M$, so
\[
\frac{m_{n+1}}{2m_n}-1=1
\qquad\text{for every }\kappa' ,
\]
which is not $O(e^{-\kappa'/2})$ as $\kappa'\to\infty$. \emph{The flat-ASCII
transcriptions that earlier drafts carried in this prose have been removed; they belong
in the all-ASCII submission sheet, where they are, and not in the typeset page.} Equivalently, at $n_{\max}=1$ the
downward form gives $m_n/(m_{n+1}/2)-1=-\tfrac12$, again a fixed nonzero constant.
\emph{Checked numerically for $\kappa'=10,20,40,80$: the relative error stays at
$1$ while $e^{-\kappa'/2}$ falls to $4\times10^{-18}$.} \textbf{The separation
condition and the one-sided cap are met with growing margin} --- $e^{-x}/\bar\delta\sim
e^{\kappa'}$ and $\kappa'/2-x\to\infty$ --- \textbf{while the admitted-error bound is met
uniformly and not with growing margin}, since $\delta_{n+1}/\bar\delta\to\tfrac14$
as $\kappa'\to\infty$; and $M\ge M_0$ is met by a fixed $M$. \emph{An earlier draft of
this note wrote that every hypothesis was met with growing margin, which is false for
those last two.} \textbf{If the implicit $O$-constant is allowed to depend
arbitrarily on $\kappa'$, the printed estimate is not quantitative enough to support
its subsequent ``harmless error'' use; under the standard uniform reading, the
family above refutes it \emph{as a general implication}, which is not the same as
refuting the conclusion for the intended operators.} What this family reaches is the
step that turns the additive estimate into a relative cascade and a uniform positive
lower bound; \textbf{the additive estimate itself is reached by the exact-trace family
of \S2}, where $\bar\delta=0$ and the remainder is genuinely present.
\emph{An earlier draft of this note said the additive estimate survived both.} \emph{An earlier draft stopped at ``the doubling
fails'', which by itself only denies exact equality --- not what the proposition
claims.}

\textbf{The finite-error spectral window is two-sided}, roughly
$\bar\delta e^{m_nM}\ll m_nM\lesssim\log(1/\bar\delta)$, together with quantitative
isolation from the next \emph{distinct} spectral value --- the third-eigenvalue term
$O(e^{-(m_n'-m_n)M}/M)$ is dropped with no hypothesis at all. For
$\bar\delta\le e^{-\kappa'}$ and $m_nM\le\kappa'/2$ the lower end requires
$m_nM\gg e^{-\kappa'/2}$, and \textbf{nothing in (H-FACT), Theorem 3.1 or Lemma 4.3
supplies that resolution floor}. \textbf{Theorem 3.1 supplies the lower bound
$m_{n_{\max}}\ge c_0$, and no hypothesis compares $c_0L_{n_{\max}}$ with that floor.}
If the certified lower bound lies below the floor, the available information does
not guarantee that the actual gap can be resolved --- \emph{the actual gap may still
be larger}. \emph{An earlier draft wrote that the smallness of $c_0$ ``can put the
gap below'' the floor, which repeats, three lines after denouncing it, the very
lower-bound-as-value slip attributed to Lemma 4.3.} \emph{This defect precedes the
$n_{\max}e^{-\kappa'/2}$ accumulation: it bites at a single step.} And the proof
inserts the error factors
``\emph{and repeats Step 4}'', so it inherits the $M$-dependence of \S2 before the
error floor $\bar\delta$ is even reached.

\paragraph{Lemma 4.3 --- a conditional identity, and an inequality used backwards.}
That $m_{n+1}=2m_n$ and $L_{n+1}=L_n/2$ imply $m_{n+1}L_{n+1}=m_nL_n$ is valid
algebra, but its premise is exactly the unestablished conclusion of Theorem 4.1.
\textbf{Lemma 4.3(a) stands only as a conditional algebraic implication}, not as an
effective invariance of the cascade.

\emph{Part (b) is worse, and it is an independent failure.} Theorem 3.1 supplies only
$m_{n_{\max}}\ge c_0$. Lemma 4.3 writes
\emph{``$m_{n_{\max}}=c_0=(\gamma N_c/8)e^{-K/\gamma}$ is exponentially small in
$1/\gamma$''} --- \textbf{it replaces a lower bound by an equality and then
uses the smallness of $c_0$ as an upper bound on the gap.} A small lower bound says
nothing about how large the gap is: $m_{n_{\max}}=1$ with $c_0=10^{-19}$ satisfies
$m_{n_{\max}}\ge c_0$ and violates the window. \textbf{The claim that the window is
automatic is not established}, and it would not be established even if Theorem 3.1's
own proof were repaired: what the window needs is an \emph{upper} bound on
$m_{n_{\max}}$, or the terminal-window condition assumed outright. The suite does not
catch this either --- checking that $c_0C_0$ is tiny checks the size of the lower
bound, not the size of the gap. \emph{An earlier draft of this note misquoted that
constant: it dropped the $\gamma$ prefactor and wrote $\kappa$ for $\gamma$ in the
exponent. Those are two different parameters --- the terminal coupling and the polymer
decay constant of Theorem 2.1(b) --- and a note whose object is the direction of an
inequality cannot make its case with a formula the paper never wrote.} \emph{An earlier draft of this note called part (b)
merely a second consumption of Theorem 3.1's numerical value. It is not: Theorem 3.1
never supplied that value, and this is a direction error of its own. It is the third
inequality-direction slip recorded in this chain.}

\paragraph{The printed control on the accumulated factorisation error ceases to be
small as $g^2\to0$.} Proposition 4.4
and the main proof produce
$m_0=(m_{n_{\max}}/2^{n_{\max}})\bigl(1+O(n_{\max}e^{-\kappa'/2})\bigr)$, and the
proof then replaces that factor by $1+o(1)$, noting that it is harmless for
$g^2\ge1/(b_0e^{\kappa'/2})$ and that ``\emph{for smaller $g^2$ the constant degrades
explicitly but remains positive}'' for $n_{\max}e^{-\kappa'/2}\le\frac12$,
``\emph{which we fold into $g_0^2$}''. \textbf{That folding runs the wrong way.}
Since $n_{\max}=\frac{1}{2b_0\log2}(1/g^2-1/\gamma)+O(1)$, at fixed $\kappa'$ one has
$n_{\max}e^{-\kappa'/2}\to\infty$ as $g^2\to0$, and the smallness condition is a
\emph{lower} bound $g^2\gtrsim c\,e^{-\kappa'/2}$, which cannot be absorbed into a
hypothesis of the form $0<g^2\le g_0^2$ --- shrinking $g_0$ makes it worse.

\textbf{What that shows, exactly.} Write the relative remainder as $R(g)$, so that the
printed estimate is an upper bound
\[
|R(g)|\ \le\ A_{\rm rel}(g)\,n_{\max}(g)\,e^{-\widetilde\kappa(g)/2}
\]
for some explicit $A_{\rm rel}$. \textbf{The printed upper bound ceases to be small} and
becomes unbounded as $g\to0$. \textbf{This does not prove that the actual error grows}:
$R\equiv0$ satisfies any such bound. What follows is only that \emph{the printed
estimate can no longer justify the replacement by $1+o(1)$} --- information insufficient,
not conclusion false. \emph{Earlier drafts of this note, and its submission sheet, said
that the accumulated error itself grows and diverges. That reads an upper bound as a
value, which is the same slip this note attributes to Lemma 4.3(b), with the inequality
turned round.}

\textbf{And the printed hypothesis does not specify coupling uniformity.} (H-FACT)
says only that $\kappa'>0$ is independent of $L_n$. It does not say that $\kappa'$ or
$C_{\rm fact}$ is independent of $g$, nor that either is uniform on $0<g^2\le g_0^2$.
\textbf{The printed (H-FACT) does not suffice under either reading}: under the
first \textbf{the printed upper bound becomes unbounded} and therefore ceases to justify
$R(g)=o(1)$ --- \textbf{this does not prove that the actual accumulated error diverges};
under the second the hypothesis is under-specified and must be supplemented.
\emph{An earlier draft of this note wrote that the accumulated error diverges here, four
lines after distinguishing a bound from a value.} Two readings, and the paper fixes neither:
\begin{itemize}\setlength{\itemsep}{0pt}
\item \textbf{if $C_{\rm fact}$, $\kappa'$ and $A_{\rm rel}$ are uniform in $g$},
the printed control parameter becomes unbounded and cannot be folded into a hypothesis
$0<g^2\le g_0^2$;
\item if they may depend on $g$, write
$\widetilde\kappa(g)=\kappa'(g)-\log C_{\rm fact}(g)$, and \textbf{choose an explicit
$A_{\rm rel}(g)\ge1$} with
$|R(g)|\le A_{\rm rel}(g)\,n_{\max}(g)\,e^{-\widetilde\kappa(g)/2}$ --- a specified
function, not ``the'' implicit constant of an $O(\cdot)$, which is not unique and can be
scaled up at will; $A_{\rm rel}\ge1$ costs nothing, since one may replace any bound by
its maximum with $1$, and it keeps $\log A_{\rm rel}$ non-negative. The growth condition
is then imposed on that specified bound:
\[
\widetilde\kappa(g)-2\log n_{\max}(g)-2\log A_{\rm rel}(g)\longrightarrow+\infty ,
\]
If $A_{\rm rel}$ is only known to be bounded above, \textbf{a sufficient
condition} --- not an equivalent one --- is $\widetilde\kappa(g)=2\log
n_{\max}(g)+\omega(1)$, which since $n_{\max}\asymp g^{-2}$ reads
$\widetilde\kappa(g)\ge4\log(1/g)+\omega(1)$; \textbf{the chosen explicit bound
$A_{\rm rel}$ must be uniformly bounded in $g$} for that reduction to be available at
all. For a
factor merely bounded by some $\eta<1$ rather than $o(1)$, the requirement relaxes to
$\widetilde\kappa\ge2\log n_{\max}+2\log(A_{\rm rel}/\eta)$. \textbf{None of this is an
$O(1)$ statement}: a bounded difference keeps the factor below a constant, it does not
send it to zero. \textbf{That growth condition is an additional missing hypothesis.}
\emph{An earlier draft of this note wrote $+O(1)$ where $\omega(1)$ is needed, and left
$A_{\rm rel}$ out of the condition while admitting three lines later that it too needs
uniform control.}
\end{itemize}
\emph{An earlier draft of this note said the flat verdict was ``the first reading
only''. That was itself wrong: the printed hypothesis fails under \emph{both}. The fault
of the flat verdict was that it did not say which reading it was arguing from, not that
its verdict was too strong.} \textbf{This propagates to Theorem 1.1's own constant.} The paper declares
$c(N_c)>0$ to depend only on $N_c$ and the universal constants of Balaban's
construction; that cannot be maintained unless $C_{\rm fact}$, $\kappa'$ and the
implicit $O$-constants are themselves controlled uniformly over the admitted data.
What would suffice: a two-sided coupling window rather than the half-line;
per-scale errors summable uniformly in $n_{\max}$; a bound improving with $g$ or with
the scale; or a proof that the accumulation renormalises without the $n_{\max}$ factor.

\paragraph{Remark 4.5 --- half stands.} That the temporal nonlocality of Balaban's
effective actions ``\emph{is not an artifact}'' stands. That (H-FACT) is
``\emph{the sharply isolated open point of this paper}'' does \textbf{not}: there are
now at least three separate unbuilt bridges --- the Euclidean-action/ground-state
measure identification (\S1.1), the four-dimensional-action/three-dimensional-operator
construction (\S1.3), and the fixed-operator gap extraction (\S2).

\paragraph{Theorem 1.1 --- not established at several distinct points audited
below, and its printed quantifier is not justified either.} \textbf{The terminal-gap
derivation and the automatic-window claim are separate failures}, not one input
consumed twice; it now also loses the transport of the gap from the terminal scale
back to the original one. \textbf{The automatic-window failure is independent of the
terminal-gap proof}: it
reverses the role of a lower bound and would survive any repair of Theorem 3.1
(\S3). The failure of the transport step is a third, independent defect.
\emph{An earlier draft of this note called the first two ``two consumptions of the
same broken input''. That was still true of the terminal gap and false of the
window, and it survived into the same round that proved it false.}

\emph{And a fourth debt, of domain rather than of proof.} The blocking sequence
$\Lambda_k=(\mathbb Z/(L/2^k)\mathbb Z)^4$ requires \textbf{admissible volumes}:
Remark 2.2 and Step 2 restrict to \textbf{$2^{n_{\max}+1}\le L\le C_0\,2^{n_{\max}}$},
equivalently $L=2^{n_{\max}}q$ with $2\le q\le C_0$ --- so $L\in2^{n_{\max}}\mathbb N$
alone is not enough: it admits $L=2^{n_{\max}}$, a terminal slice of extent one, which
Remark 2.2 itself excludes. \textbf{And the lower endpoint depends on $g$} through
$n_{\max}$, so it cannot be replaced by a fixed $L_0$ as $g^2\to0$. Theorem 1.1 nevertheless quantifies
over every $L$ in an interval, and no rounding, interpolation or irregular-blocking
argument is given for the rest. \textbf{The statement is not literally established
with its printed quantifier over all $L$.} \emph{This is the same correction the
companion erratum applies to Theorem 6.3 of \texttt{ai.viXra:2602.0035}; it had not
been propagated back to the paper the cascade comes from.}

\section{Two residual pieces of structure, with their exact limitations}

Neither is a printed statement that survives as printed. Both are pieces of structure
that a reformulated route could reuse, and each is stated here with what it still owes.

\subsection{The transfer operator is on the intended sector, but the measure is unstated}

Theorem 4.1 writes $\Hil_n=L^2(B^{\rm spat}_n)$ without saying whether the measure is
$\nu_{\rm spat}=\pi_\#\mathrm{Haar}$, the Riemannian volume of the regular stratum, or
another. \emph{An earlier draft of this note called this ``the right object, by
(I1)''; that was premature and is withdrawn}, and it contradicted the appended page,
which records that this paper's spatial domain is not yet fixed. The accurate
statement is: \textbf{Theorem 4.1 intends a transfer operator on the gauge-invariant
spatial sector, but the printed notation does not specify the measure; identification
(I1) applies only after $\nu_{\rm spat}=\pi_\#\mathrm{Haar}$ is taken explicitly.}

\subsection{The terminal oscillation bound needs a two-sided coupling hypothesis}

Step 2 bounds $\operatorname{osc}(V_W)=O(1)$ by the plaquette count and
$\sum_X|\epsilon_{n_{\max}}(X)|\le C_\epsilon(\kappa,d)<\infty$, both correctly. It
then substitutes $g^2_{n_{\max}}=\gamma$ --- equivalently
$\beta_{n_{\max}}=2N_c/\gamma$ --- and concludes $\operatorname{osc}(f)\le K/\gamma$.

But Theorem 2.1(c) supplies only $g^2_{n_{\max}}\le\gamma$, i.e.\
$\beta_{n_{\max}}\ge 2N_c/\gamma$ --- a \emph{lower} bound on $\beta_{n_{\max}}$. An
upper bound on $\beta_{n_{\max}}\operatorname{osc}(V_W)$ requires an \emph{upper} bound
on $\beta_{n_{\max}}$, and the printed inequality runs the other way. \emph{An earlier
draft of this note asserted that this bound ``does establish what it claims
uniformly''; that is withdrawn.}

\emph{And it is used twice, not once.} The last line of the same proof concludes
\[
m_{n_{\max}}=\tfrac{1}{2}g^2_{n_{\max}}\lambda_1
\;\ge\;\tfrac{1}{8}\gamma N_c\,e^{-K/\gamma},
\]
replacing the prefactor $g^2_{n_{\max}}$ by $\gamma$ inside a \emph{lower} bound.
The hypothesis $g^2_{n_{\max}}\le\gamma$ does not license that either: it makes the
prefactor smaller, not larger. So the unsupported replacement
$g^2_{n_{\max}}=\gamma$ occurs at both ends of the proof --- once to bound
$\beta_{n_{\max}}$ from above in the oscillation estimate, once to bound the kinetic
prefactor from below --- and the printed one-sided hypothesis supplies neither
inference. \textbf{A two-sided terminal-coupling bound is required} ---
$c\gamma\le g^2_{n_{\max}}\le\gamma$ for some $c>0$, or directly
$\beta_{n_{\max}}\le C_\beta/\gamma$ --- and it repairs both places at once, with constants
degraded by $c$. (The direction of this slip is the same as one corrected in
\texttt{ai.viXra:2602.0052}: substituting an endpoint into a bound that moves the
wrong way.)

What survives is only the finite-dimensional oscillation estimate for the explicit
Euclidean effective action, namely
\[
\operatorname{osc}(f)\;\le\;\beta_{n_{\max}}C_W+2C_\epsilon ,
\qquad C_W:=\operatorname{osc}(V_W),\quad
C_\epsilon:=C_\epsilon(\kappa,d).
\]
This is an explicit finite-dimensional inequality. It becomes uniform in the claimed
parameters only after a two-sided terminal-coupling bound is supplied --- note that
the right-hand side still contains $\beta_{n_{\max}}$, which is exactly the quantity
not yet bounded above. And \textbf{it does not control the ground-state exponent}
$f_0=-\log\psi_0^2$; it therefore neither supplies an $(\text{A-osc})$-type hypothesis
about the ground state nor removes the separate identification required by (I3).

\section{The citation loop, and what a clean repair would need}

Theorem 3.1 proves the terminal gap by ``\emph{applying the argument of
{\rm[3, Theorem 5.3]}}'', where [3] is \texttt{ai.viXra:2602.0035}; that paper cites
this one as discharging its (H1). The loop is not logically vicious --- this paper
re-runs the method with its own inputs --- but \textbf{this paper is not an independent
confirmation of that one}, and neither discharges a hypothesis of the other.

A standalone terminal lemma stated on $(B,\nu,K_\nu)$ would \emph{not} by itself
suffice: for the Euclidean measure $e^{-S_{\rm eff}}\nu$ it would remain a statement
about a Euclidean diffusion. To discharge a Hamiltonian gap the lemma must be stated
either for $\mu_0=\psi_0^2\nu$, or together with an explicit transfer-operator
construction identifying the two quantities.

\section{Scope}

Not examined here: Balaban's package, the numerical suite beyond the observation of
\S2, Theorem 2.1 and the coupling-evolution erratum of Remark 2.3, and the existence
of a positive self-adjoint trace-class transfer matrix at each scale, which Theorem
4.1 assumes as part of its hypothesis. \emph{This audit does not refute} (H-BAL)
\emph{or} (H-FACT) --- both are, as printed, unusually candid --- but those hypotheses
do not suffice for the printed conclusions without further inputs, \textbf{including}
all of the following:
\begin{enumerate}\setlength{\itemsep}{0pt}
\item a \emph{globally specified} Euclidean measure with a closed form, along one of
the two branches of (I2b) --- and hence either control of $\operatorname{osc}(\log w)$
on route (A), or a theorem identifying the $\nu$-based measure as the correct object
on route (B) (\S1.1--1.2);
\item a four-dimensional-action to three-dimensional-space transfer-operator or
Hamiltonian construction (\S1.3);
\item the ground-state-measure identification (I3);
\item the renormalisation $\kappa'\mapsto\kappa'-\log C_{\rm fact}$ that
Proposition 4.4 needs in order to consume (H-FACT) as printed when $C_{\rm fact}>1$,
together with $\kappa'>\log C_{\rm fact}$ (\S3);
\item \emph{either} a restriction of the theorem to a coupling window chosen so
that $A_{\rm rel}\,n_{\max}(g)\,e^{-\widetilde\kappa/2}\le\eta<1$ on it, if
$C_{\rm fact}$, $\kappa'$ and $A_{\rm rel}$ are uniform in $g$ --- a compact window
alone only makes $n_{\max}$ finite, not small against the error floor, and in any case
this changes the statement rather than repairing its printed quantifier over all
sufficiently small couplings --- \emph{or}, if they may depend on $g$, the growth
condition
$\widetilde\kappa(g)-2\log n_{\max}(g)-2\log A_{\rm rel}(g)\to+\infty$ (\S3);
\item a two-sided terminal-coupling bound (\S4.2);
\item the admissible-volume condition $2^{n_{\max}+1}\le L\le C_0\,2^{n_{\max}}$ for
the blocking sequence, whose lower endpoint moves with $g$ (\S3);
\item an \emph{upper} bound on $m_{n_{\max}}$, or the terminal-window condition as a
hypothesis, for Lemma 4.3(b) (\S3);
\item per-scale factorisation errors that stay summable as $g^2\to0$, which the
printed (H-FACT) does not guarantee under either reading of its
coupling-uniformity (\S3);
\item {\sloppy a two-sided spectral-resolution window for the approximate
trace extraction, including $\bar\delta e^{m_nL_n}\ll m_nL_n$ and
quantitative isolation from the next distinct spectral value (\S3);\par}
\item for the cascade (\S2), \emph{either} a fixed-spatial-slice family
\emph{together with} an identity valid for all temporal extents over that fixed space,
\emph{or} an operator relation, \emph{or} a \emph{normalised} spectral-convergence
 theorem in the sense of \S2, with normalisation relative to the actual finite-$M$
 first-excited contributions, \emph{subexponential control of the excited
 multiplicities, and uniform control of the total excited trace tail}.  The displayed
 multiplicity and tail conditions are necessary but not sufficient without that
 finite-$M$ normalisation; ordinary spectral convergence and the old limiting-scale
 normalisation are insufficient, positivity improving removes only the
 top-eigenvalue multiplicity, and the fixed slice alone does not suffice.
\end{enumerate}
\emph{An earlier draft of this section listed only items 3--5. Items 1 and 2 are
defects this same note had already separated, and the closing list dropped them; it is
written with ``including'' now because a closing list of obligations that reads as
exhaustive is the same failure as a survivors list closed too early.}

\vspace{4pt}
\noindent\rule{\textwidth}{0.8pt}

\noindent\footnotesize
\textbf{Provenance.} \S1 was found by asking which measure each paper of this chain
uses, and by auditing the cited Ricci theorem rather than the citing sentence.
\textbf{\S2, \S4.1 and \S4.2 came from external audits of the earlier drafts of this
very note}, which had certified the doubling cascade as surviving intact, the spatial
sector as the right object, and the oscillation bound as uniform. All three were
premature. A second audit round supplied the rest: the \emph{second} inverted use of
the coupling hypothesis in the last line of the same proof, the positivity condition
$a,b>0$ without which the abstract example has the wrong top eigenvalue, the
multiplicity defect of Step 3, and the distinction between three proof failures and
three independent routes. A third round found the drift \emph{on the appended page
itself}, which had recorded this paper's printed measure against $\nu$ instead of
$\mathrm{Vol}_B$ --- the very substitution its new (I2b) forbids --- and then found
that first correction overshooting, asserting the two measures \emph{are} different
when that needs $w$ non-constant and is not proved here. A fourth round supplied the
three-eigenvalue family, the thresholds $W(1)$ and $W(2)$ for the window, and the
separation between a false general implication and a refuted printed conclusion ---
the note had crossed it in the direction that overstates. The pattern --- a survivors list closed by checking what a
statement depends on rather than by reading its proof --- is the same one this
correction campaign has now recorded five times, and is the reason the appended page
exists. That page is reproduced verbatim: paraphrase is how the columns drift.

\end{document}
